\chapter{Vorprojekt: KI-basierte Roboterbahnen}
\label{ch:vorprojekt}

Vor dem vorliegenden Projekt hat die Projektgruppe an einer Aufgabenstellung in der Industrie
gearbeitet. Bei der GROB-WERKE GmbH \& Co.\ KG sollte das Thema \enquote{Erstellung von
KI-basierten Roboterpfaden für die Robotersimulation und Inbetriebnahme} bearbeitet werden. Nach
einer Einarbeitung von acht Tagen vor Ort wurde das Projekt bearbeitet. Während der Bearbeitung
wurde die Zusammenarbeit beendet, und die Gruppe hat sich die Aufgabe dieser Arbeit selbst
gestellt (\cref{ch:einleitung}). Die Ergebnisse des Vorprojekts sind deshalb begrenzt.

Dieses Kapitel fasst zusammen, was bis dahin erarbeitet wurde. Es beschreibt die Aufgabe, die
Werkzeuge und das Vorgehen bei der Bahnplanung in der Robotersimulation und die ersten
Erfahrungen mit Werkzeugen zur automatischen Bahnplanung. Am Ende wird eingeordnet, welche
Erkenntnisse für den Demonstrator relevant sind oder eine mögliche Erweiterung wären.

% ------------------------------------------------------------------------------
\section{Aufgabenstellung und Werkzeugumgebung}
\label{sec:vp-aufgabe}

Bevor eine Roboterzelle aufgebaut wird, werden ihre Bewegungsabläufe in einer
Robotersimulation geplant und geprüft. Dafür wird ein digitales Modell der Zelle aus den
Konstruktionsdaten verwendet, mit Robotern, Greifern, Vorrichtungen und Schutzeinrichtungen.
Die dort erstellten Bahnen werden später auf die realen Roboter übertragen. Je genauer die
Bahnen in der Simulation stimmen, desto weniger muss bei der Inbetriebnahme an der Anlage
nachgebessert werden.

Bei GROB wird dafür Siemens Process Simulate eingesetzt. Die Bahnen werden dort bisher von Hand
erstellt. Das kostet bei großen Zellen mit mehreren Robotern viel Zeit. Ziel des Vorprojekts war
es zu untersuchen, ob sich Bahnen mit KI-basierten Werkzeugen automatisch erzeugen lassen und
ob diese Bahnen den Anforderungen von GROB genügen. Ergebnis sollte in jedem Fall eine
kollisionsfreie und auf die Taktzeit optimierte Roboterbahn sein.

Die Aufgabe wurde auf einen Teil einer bestehenden Zelle eingegrenzt. Betrachtet wurde ein
Handhabungsprozess, bei dem ein Roboter Teile zwischen festen Positionen umsetzt. Prozesse wie
das Kleben wurden nicht betrachtet. Sie sind deutlich komplexer und für die Fragestellung
nicht notwendig. Verwendet
wurde ein KUKA-Roboter, weil sich dieser gut in Process Simulate abbilden lässt und einfacher zu
handhaben ist.

In der Einarbeitung wurden vier Themen behandelt: Kollisionserkennung, Roboter- und
Zellsicherheit, das manuelle Erstellen von Bahnen und erste Werkzeuge zur automatischen
Bahnplanung. Sie werden in den folgenden Abschnitten beschrieben, bevor
\cref{sec:vp-vorgehen} das Vorgehen im Projekt zeigt.

\todo[inline]{Quelle für Siemens Process Simulate (Herstellerdokumentation) ergänzen, falls
gewünscht.}

% ------------------------------------------------------------------------------
\section{Kollisionserkennung}
\label{sec:vp-kollision}

Die Simulation prüft nicht automatisch jedes Bauteil der Zelle gegen jedes andere. Stattdessen
legt der Anwender \emph{Kollisionssets} an. Ein Kollisionsset besteht aus zwei Listen von
Objekten: Die Objekte der ersten Liste werden gegen die Objekte der zweiten Liste geprüft
(\cref{fig:vp-kollisionsset-editor}). Typischerweise stehen in der ersten Liste der Roboter und
sein Greifer. In der zweiten Liste stehen die Vorrichtungen, Werkstücke und Anlagenteile, mit
denen er zusammenstoßen könnte. Paare, die sich ohnehin berühren, etwa der Greifer und der
Roboterflansch, werden so von der Prüfung ausgenommen. Das spart Rechenzeit und vermeidet
Fehlmeldungen.

\begin{figure}[htbp]
  \centering
  \includegraphics[width=0.6\textwidth]{grob/collision_set_editor.png}
  \caption{Anlegen eines Kollisionssets: Die Objekte links werden gegen die Objekte rechts
  geprüft}
  \label{fig:vp-kollisionsset-editor}
\end{figure}

Für jedes Kollisionsset lassen sich zwei Grenzwerte einstellen. Der \emph{Beinahe-Abstand}
(\enquote{Near Miss}) meldet, wenn sich zwei Objekte näher als ein festgelegter Abstand kommen,
ohne sich zu berühren. Die \emph{zulässige Durchdringung} erlaubt geringe Überschneidungen, zum
Beispiel wenn ein Greifer ein Werkstück greift. Während der Simulation hebt die Software die
Objekte eines Sets farbig hervor und listet die kollidierenden Paare auf
(\cref{fig:vp-kollisionsset-anzeige}).

\begin{figure}[htbp]
  \centering
  \includegraphics[width=0.9\textwidth]{grob/collision_set_anzeige.png}
  \caption{Hervorgehobene Objekte eines Kollisionssets in einer Roboterzelle}
  \label{fig:vp-kollisionsset-anzeige}
\end{figure}

Für die automatische Bahnplanung sind die Kollisionssets die Grundlage. Ein Planungswerkzeug
kann nur solche Kollisionen vermeiden, die es auch prüft. Fehlt ein Objekt im Set, kann eine
Bahn in der Simulation kollisionsfrei erscheinen und an der realen Anlage trotzdem zu einem
Zusammenstoß führen. Jedes Werkzeug zur automatischen Bahnplanung muss die Kollisionssets
deshalb berücksichtigen.

% ------------------------------------------------------------------------------
\section{Roboter- und Zellsicherheit}
\label{sec:vp-sicherheit}

Neben der Kollisionsprüfung in der Simulation überwacht die Sicherheitssteuerung des realen
Roboters, dass er bestimmte Bereiche nicht verlässt oder nicht betritt. Diese Bereiche werden
ebenfalls in der Simulation festgelegt und dort sichtbar gemacht. In der Einarbeitung wurde
dafür eine Erweiterung von Process Simulate für die Sicherheitskonfiguration der eingesetzten
KUKA-Roboter verwendet. Drei Elemente wurden behandelt.

\paragraph{Zellbereich.}
Der Zellbereich beschreibt den Raum, den der Roboter insgesamt nicht verlassen darf. Er wird als
Prisma angegeben: eine Grundfläche aus bis zu zehn Eckpunkten in der Ebene und eine untere und
obere Höhengrenze (\cref{fig:vp-zellbereich}).

\begin{figure}[htbp]
  \centering
  \includegraphics[width=0.75\textwidth]{grob/cell_space.png}
  \caption{Festlegen des Zellbereichs über Eckpunkte der Grundfläche und Höhengrenzen}
  \label{fig:vp-zellbereich}
\end{figure}

\paragraph{Überwachungsräume.}
Innerhalb des Zellbereichs lassen sich weitere Räume als Quader festlegen
(\cref{fig:vp-ueberwachungsraum}). Ein Quader wird durch ein Bezugskoordinatensystem und die
minimalen und maximalen Abstände in jeder Achse beschrieben. Für jeden Raum wird eingestellt, ob
der Roboter ihn nicht verlassen darf (Arbeitsraum) oder nicht betreten darf (Schutzraum), ob er
immer oder nur unter Bedingungen aktiv ist und ob der Roboter an der Grenze anhält. Ein
Schutzraum wird zum Beispiel um eine Vorrichtung gelegt, in die der Roboter nicht eingreifen
soll (\cref{fig:vp-ueberwachungsraum-anzeige}).

\begin{figure}[htbp]
  \centering
  \includegraphics[width=0.85\textwidth]{grob/monitoring_space.png}
  \caption{Parameter eines Überwachungsraums vom Typ Schutzraum}
  \label{fig:vp-ueberwachungsraum}
\end{figure}

\begin{figure}[htbp]
  \centering
  \includegraphics[width=0.5\textwidth]{grob/monitoring_space_anzeige.png}
  \caption{Darstellung eines Schutzraums (transparenter Quader) um eine Vorrichtung}
  \label{fig:vp-ueberwachungsraum-anzeige}
\end{figure}

\paragraph{Werkzeugkugeln.}
Damit die Sicherheitssteuerung nicht nur den \ac{tcp}, sondern das ganze Werkzeug überwachen
kann, wird der Greifer durch Kugeln angenähert (\cref{fig:vp-werkzeugkugeln}). Jede Kugel ist
durch ihren Mittelpunkt relativ zum Roboterflansch und ihren Radius festgelegt. Die Kugeln
müssen den Greifer vollständig umschließen. Zu große Kugeln schränken den nutzbaren
Arbeitsraum dagegen unnötig ein. Die Steuerung prüft dann für jede Kugel, ob sie einen
Überwachungsraum verletzt. Kugeln sind dafür gut geeignet, weil sich der Abstand einer Kugel zu
einer Fläche mit wenig Rechenaufwand bestimmen lässt.

\begin{figure}[htbp]
  \centering
  \includegraphics[width=0.85\textwidth]{grob/tool_spheres.png}
  \caption{Festlegen der Werkzeugkugeln über Mittelpunkt und Radius}
  \label{fig:vp-werkzeugkugeln}
\end{figure}

\begin{figure}[htbp]
  \centering
  \includegraphics[width=0.8\textwidth]{grob/tool_spheres_anzeige.png}
  \caption{Greifer mit den umschließenden Werkzeugkugeln}
  \label{fig:vp-werkzeugkugeln-anzeige}
\end{figure}

% ------------------------------------------------------------------------------
\section{Manuelle Bahnplanung}
\label{sec:vp-manuell}

Bei GROB werden die Bahnen nach einem festen Vorgehen erstellt. Gegeben sind die Startpose und
die Zielpose, zum Beispiel eine Entnahme- und eine Ablageposition. Zu beiden Posen wird
zuerst eine Vorpose erzeugt. Von der Vorpose aus fährt der Roboter mit einer \ac{lin} in die
Start- bzw.\ Zielpose, wie es auch in \cref{sec:robotik} beschrieben ist.

Die Verbindung zwischen den beiden Vorposen wird schrittweise aufgebaut
(\cref{fig:vp-manuell}). Der Roboter wird in der Simulation in Gelenkkoordinaten von der einen
Vorpose zur anderen verfahren, bis eine Kollision auftritt. An dieser Stelle wird er aus der
Kollision herausbewegt, möglichst durch das Verfahren nur einer Achse. Die so gefundene Stellung
wird als Zwischenpunkt gespeichert, und die Bahn wird von dort aus erneut betrachtet. Das wird
wiederholt, bis die gesamte Bahn kollisionsfrei ist.

\begin{figure}[htbp]
  \centering
  \begin{tikzpicture}[pap]
    \node[pap term] (start) {Startpose $S$, Zielpose $Z$};
    \node[pap proc, below=of start] (vor) {Vorposen $V_S$ und $V_Z$ erzeugen};
    \node[pap proc, below=of vor] (lin1) {\acs{lin} von $S$ nach $V_S$};
    \node[pap proc, below=of lin1] (jog) {In Gelenkkoordinaten\\Richtung $V_Z$ verfahren};
    \node[pap dec, below=10mm of jog] (koll) {Kollision?};
    \node[pap proc=56mm, right=22mm of koll] (raus) {Aus der Kollision herausfahren, möglichst mit nur einer Achse};
    \node[pap proc=56mm] (zp) at (raus |- jog) {Stellung als Zwischenpunkt\\speichern (\acs{ptp})};
    \node[pap proc, below=of koll] (ptp) {\acs{ptp} nach $V_Z$};
    \node[pap proc, below=of ptp] (lin2) {\acs{lin} von $V_Z$ nach $Z$};
    \node[pap term, below=of lin2] (ende) {kollisionsfreie Bahn};
    \draw[pap arr] (start) -- (vor);
    \draw[pap arr] (vor) -- (lin1);
    \draw[pap arr] (lin1) -- (jog);
    \draw[pap arr] (jog) -- (koll);
    \draw[pap arr] (koll) -- node[pap lbl] {ja} (raus);
    \draw[pap arr] (raus) -- (zp);
    \draw[pap arr] (zp) -- (jog);
    \draw[pap arr] (koll) -- node[pap lbl, right=1pt] {nein, $V_Z$ erreicht} (ptp);
    \draw[pap arr] (ptp) -- (lin2);
    \draw[pap arr] (lin2) -- (ende);
  \end{tikzpicture}
  \caption{Manuelle Bahnplanung zwischen zwei Vorposen}
  \label{fig:vp-manuell}
\end{figure}

Ziel dieses Vorgehens ist, die Bewegungen der einzelnen Achsen klein zu halten. Oft wurde
außerdem gewünscht, dass die Achsen~4 und~6 bei \SI{0}{\degree} bleiben. Das ist nicht
zwingend, vermeidet aber unnötige Drehungen der Handachsen und damit eine starke Belastung der
Energiezuführung zum Greifer. Außerdem bleiben die Bahnen für den Inbetriebnehmer
nachvollziehbar.

% ------------------------------------------------------------------------------
\section{Vorgehen im Projekt}
\label{sec:vp-vorgehen}

Die betrachtete Zelle ist bereits seit einigen Jahren umgesetzt. Deshalb wurde zuerst das
bestehende Projekt analysiert. Die Taktzeiten der dort umgesetzten Bahnen dienen als
Vergleichsbasis für alle neu erzeugten Bahnen.

Anschließend wurde ein leeres Bahnmodul angelegt. Es enthält nur die Start- und Zielposen und
die zugehörigen Vorposen, aber keine Zwischenpunkte. Alle Werkzeuge zur automatischen
Bahnplanung erhalten dasselbe Bahnmodul als Ausgangspunkt. Dadurch sind ihre Ergebnisse
untereinander und mit der bestehenden Bahn vergleichbar.

Zusätzlich hat die Projektgruppe selbst Bahnen nach dem manuellen Vorgehen aus
\cref{sec:vp-manuell} erstellt. Diese Bahnen zeigen, welches Ergebnis mit wenig Erfahrung in
der Bahnplanung erreicht wird, und dienen als weiterer Vergleich. Gleichzeitig hat die Gruppe
dabei das Programm und die Bahnplanung selbst kennengelernt.

Damit ergeben sich drei Vergleichswerte für eine automatisch erzeugte Bahn: die bestehende,
von erfahrenen Planern erstellte Bahn, die eigene manuelle Bahn und die Bahnen der
anderen Werkzeuge.

\todo[inline]{Falls vorhanden: Taktzeiten der bestehenden Bahn, der eigenen manuellen Bahn und
der Bahn aus dem Robotic Automatic Path Planner als Tabelle ergänzen.}

% ------------------------------------------------------------------------------
\section{Werkzeuge zur automatischen Bahnplanung}
\label{sec:vp-werkzeuge}

Für die automatische Erzeugung der Bahnen waren zwei Werkzeuge vorgesehen. Als dritte
Möglichkeit sollten Verfahren betrachtet werden, die nicht in Process Simulate enthalten sind.

\paragraph{Robotic Automatic Path Planner.}
Der Robotic Automatic Path Planner ist in Process Simulate integriert. Er erzeugt aus einer
vorhandenen Bahn mit Start- und Zielpose eine kollisionsfreie Bahn und berücksichtigt dabei die
Kollisionssets. Der Anwender kann zwischen mehreren Optimierungszielen wählen, zum Beispiel
einer möglichst kurzen Taktzeit.

Das Werkzeug wurde auf das leere Bahnmodul angewendet, und Bahnen konnten erfolgreich
automatisch erstellt werden. Bei der Optimierung auf die Taktzeit fiel jedoch auf, dass das Werkzeug auch
die programmierten Geschwindigkeiten ändert. Betroffen war zum Beispiel die geradlinige Bewegung
von der Vorpose in die Zielpose. Diese Bewegung wird bewusst langsam programmiert, damit das
Werkstück sicher abgelegt oder gegriffen wird. Wird sie beschleunigt, kann das Werkstück
verrutschen, oder Vorrichtung und Greifer werden stärker belastet. Eine automatisch erzeugte
Bahn muss deshalb darauf geprüft werden, ob solche festen Vorgaben erhalten bleiben.

\todo[inline]{Weitere Optimierungsziele des Robotic Automatic Path Planner nennen (z.\,B. Abstand
zu Hindernissen, Achsbewegung), falls bekannt.}

\paragraph{Realtime Robotics.}
Als zweites Werkzeug war eine Lösung der Firma Realtime Robotics vorgesehen. Die Lizenz stand
bis zum Ende der Zusammenarbeit nicht zur Verfügung, deshalb konnte das Werkzeug nicht getestet
werden.

\paragraph{Alternative Verfahren.}
\todo[inline]{Kommilitone: Alternative Verfahren zur Bahnplanung außerhalb von Process Simulate
beschreiben (Ansätze, betrachtete Werkzeuge, Einschätzung). Neue Quellen nur mit geprüfter DOI.}

% ------------------------------------------------------------------------------
\section{Bedeutung für das vorliegende Projekt}
\label{sec:vp-bedeutung}

Die Aufgabe des vorliegenden Projekts unterscheidet sich vom Vorprojekt. Im Mittelpunkt steht
nicht die Planung von Bahnen in der Simulation, sondern das Erkennen, Greifen und Prüfen eines
Bauteils am realen Roboter. Mehrere Erkenntnisse aus dem Vorprojekt sind trotzdem übertragbar:

\begin{description}
  \item[Vorpose und geradlinige Anfahrt] Die Greif- und Ablegesequenz des Demonstrators folgt
    demselben Muster wie bei GROB: Eine Vorpose wird mit einer Gelenkbewegung angefahren, das
    Ziel dann langsam und geradlinig (\cref{sec:m2-sequenz}). Die Erfahrung mit dem
    Robotic Automatic Path Planner zeigt, dass diese langsame Anfahrt als feste Vorgabe gelten
    muss, auch wenn Bahnen später automatisch erzeugt oder optimiert werden.
  \item[Sichere Bereiche] Die Robotersteuerung des Demonstrators prüft vor jeder Bewegung,
    ob Zielpose und Vorposition in einem festgelegten Quader liegen (\cref{sec:m2-sicherheit}).
    Das entspricht in vereinfachter Form einem Arbeitsraum. Geprüft wird dabei nur der \ac{tcp}.
    Eine mögliche Erweiterung wäre, Greifer und Bauteil wie bei den Werkzeugkugeln durch Kugeln
    anzunähern und zusätzlich Schutzräume um Kameras und Prüfstation festzulegen.
  \item[Kollisionsprüfung in der Simulation] Im Demonstrator werden die Posen geteacht, und die
    Wege dazwischen führen durch freien Raum über dem Tisch. Eine Kollisionsprüfung ist deshalb
    nicht nötig. Würde die Zelle erweitert, könnten die Bahnen vorab in einer Simulation mit
    einem Zellmodell und Kollisionssets geprüft werden. Für den eingesetzten Roboter bietet der
    Hersteller dafür eine Offline-Simulation an~\cite{neura2024offlinesim}.
  \item[Randbedingungen für KI-Verfahren] Im Vorprojekt sollte eine KI die Bahnen erzeugen, im
    Demonstrator bestimmt ein lernbasiertes Verfahren optional die Greifpose
    (\cref{ch:kigreifen}). In beiden Fällen gilt: Das Ergebnis der KI muss die
    Randbedingungen der Anlage einhalten, also kollisionsfrei sein, in den erlaubten Bereichen
    liegen und feste Vorgaben wie Geschwindigkeiten beibehalten. Im Demonstrator prüft deshalb
    die Robotersteuerung jede Zielpose gegen die Arbeitsraumgrenzen, unabhängig davon, welches
    Modul sie liefert.
\end{description}
